\section{Scope, Chronology, and Evidence}

\subsection{Subject, names, and range}

Subject: Saint Ambrose of Milan (Aurelius Ambrosius), Roman advocate, consular governor (\latin{consularis}) of Aemilia and Liguria, bishop of Milan (374--397), Father and Doctor of the Church. Whole-life range: born in Gaul, probably at Trier (Augusta Treverorum, modern Trier, Germany), c.~339--340 (older alternative c.~333); died at Milan (Mediolanum, modern Milan, Italy) on Holy Saturday, 4~April 397. Principal historical places: Trier, Rome, Sirmium, Milan, Aquileia, Bononia (Bologna), Florence, and Vercellae (Vercelli). Ecclesial state of life: secular Roman aristocrat and catechumen until his sudden popular election in 374; baptized Catholic, ordained, and consecrated bishop of Milan within eight days; metropolitan pastor and celibate diocesan bishop. Liturgical commemoration: 7~December (anniversary of episcopal ordination) in the Ambrosian and Roman Rites; 4~April (deposition) in the \work{Martyrologium Romanum}.

\subsection{Audience, scope, and boundary}

Intended for a serious general reader, student, or researcher seeking one continuous, source-controlled historical life of Ambrose and a clear account of his ecclesial reception. The governing questions open the first section. The historical synthesis---that a Roman magistrate suddenly made bishop reshaped Latin Nicene orthodoxy, sacramental catechesis, clerical ethics, Western hymnody, and Church-state accountability, while leaving both heroic witnesses and troubling episodes such as Callinicum---is historical and analytical rather than ecclesiastical. No doctrinal adjudication, forensic authentication of relics, miracle certification, independent patristic specialist review, canonical opinion, rights opinion, or ecclesiastical approval is claimed.

\subsection{Corpus, editions, and reporting ceilings}

All quoted English texts of Ambrose's treatises and letters come directly from the unrevised public-domain CCEL printing of \work{A Select Library of Nicene and Post-Nicene Fathers}, Series~2, Vol.~10, translated by H.~de Romestin, E.~de Romestin, and H.~T.~F. Duckworth (Edinburgh/New York, 1896), including Symmachus's \work{Relatio} III and the \work{Sermon against Auxentius}. \work{De Excessu Fratris Satyri} is bound through its constituent segment within that same volume. Latin wording of the \work{Hexameron} comes directly from the tracked Migne \work{Patrologia Latina} 14 text; verified Latin and English passages of \work{Expositio Evangelii secundum Lucam} come from K.~Schenkl's CSEL~32/4 (1902) and the 1843 Oxford \work{Catena Aurea}; verified Latin Eucharistic passages of \work{De Sacramentis} IV.5--6 come from Migne PL~16 (1880). Contemporary external testimony comes directly from tracked unrevised CCEL printings of St.~Basil's \work{Letter} CXCVII (NPNF~II.8) and St.~Augustine's \work{Confessions}, \work{Letters} 36 and 54 (NPNF~I.1), \work{The City of God} XXII.8 and \work{On Christian Doctrine} IV.21 (NPNF~I.2), and \work{On the Grace of Christ} I.43--44 (NPNF~I.5). Official liturgical reception comes directly from the 1902 \work{Martyrologium Romanum}, the eleventh-century \work{Manuale Ambrosianum} (ed.\ Magistretti, 1904), the 1640 \work{Missale Ambrosianum}, the verified 7~December proper orations of the 1962 \work{Missale Romanum} (p.~435, no.~1863) collated with the public-domain 1945 Benziger \work{New Roman Missal} (pp.~845--846), and the postconciliar General Roman Calendar inventory.

Several ancient and modern witnesses are not held as standalone full texts in the local archive and are therefore used strictly at their reporting ceiling: Paulinus of Milan's \work{Vita Ambrosii} (c.~422), Rufinus, Socrates, Sozomen, Theodoret (\work{Hist.\ Eccl.} V.18), Jerome's \work{De Viris Illustribus} and \work{Adversus Rufinum}, Ambrose's funeral orations for Valentinian~II and Theodosius~I, Boniface~VIII's 1298 decretal confirming the four Latin Doctors, modern hymnological attribution of the \work{Te Deum} to Nicetas of Remesiana, and nineteenth-century archaeological reports on the 1864 Basilica Ambrosiana excavations. Details from those works appear only as attributed secondary reports (chiefly through H.~de Romestin's 1896 Prolegomena and notes, Blomfield Jackson's NPNF~II.8 notes, and J.~G. Pilkington's NPNF~I.1 notes), with known editorial errors---such as de Romestin's misnaming of Saints Gervasius and Protasius as ``SS.\ Cosmas and Damian'' in Prolegomena~\S II---explicitly corrected against Ambrose's own \work{Letter} XXII. All source inspections and status checks carry an as-of cutoff of 27~November 2026.

\subsection{Evidence classes}

\begin{evidencekey}
\evidenceclass{A}{Subject-authored evidence: Ambrose's surviving letters, sermons, treatises, and biblical expositions (\work{De Officiis}, \work{De Excessu}, \work{De Fide}, \work{De Spiritu Sancto}, \work{De Mysteriis}, \work{De Paenitentia}, \work{De Virginibus}, \work{De Viduis}, \work{Hexameron}, \work{Expositio Lucae}, \work{Sermon against Auxentius}, \work{Epistles} XVII--XXII, XL--XLI, LI, LVII, LXI--LXIII); qualified by rhetorical occasion, polemical genre, and post-delivery revision.}
\evidenceclass{N}{Contemporary external witness: St.~Basil's \work{Letter} CXCVII (c.~375), Symmachus's \work{Relatio} III (384), St.~Augustine's \work{Confessions} V--IX, \work{Letters} 36 and 54, \work{City of God} XXII.8, \work{De Doctrina Christiana} IV.21, and \work{On the Grace of Christ} I.43--44 (preserving Pelagius's praise).}
\evidenceclass{E}{Early post-mortem tradition: Paulinus of Milan's \work{Vita Ambrosii} (c.~422), Rufinus, Socrates, Sozomen, Theodoret, and Jerome's anti-Rufinian notices; each used at its stated reporting ceiling.}
\evidenceclass{L}{Later liturgical, hagiographic, and iconographic tradition: eleventh-century \work{Manuale Ambrosianum}, the \work{Te Deum} baptismal legend, beehive iconography, the 835 porphyry sarcophagus report, and the broader ``Ambrosian'' hymnal.}
\evidenceclass{M}{Official ecclesial reception, limited to each act's actual object: 1640 \work{Missale Ambrosianum} (Cardinal Monti preface), 1902 \work{Martyrologium Romanum}, 1962 and postconciliar Roman calendars for 7~December, and Boniface~VIII's 1298 decretal at reported level.}
\evidenceclass{K}{Named modern critical and patristic reference: Maurist/Benedictine editorial apparatus, K.~Schenkl's CSEL~32/4 (1902), H.~de Romestin's NPNF~II.10 Prolegomena (1896), B.~Jackson's NPNF~II.8 notes (citing H.~B. Swete), and modern chronological/hymnological consensus at reported ceiling.}
\evidenceclass{S}{Source-grounded historical synthesis, distinct from individual witnesses.}
\end{evidencekey}

\subsection{Chronology}

Dates are historical claims with stated bases and competing alternatives preserved. Where fourth-century letters lack consular year dates, sequences follow Ambrose's internal historical cross-references and the editorial chronology of NPNF~II.10 checked against modern chronological alternatives.

\begin{lifetimeline}
\lifeevent{c.~339--340 (alt.\ c.~333)}{Birth in Gaul, probably at Trier, youngest child of praetorian prefect Ambrosius}{K (\work{Prolegomena}~\S III, \S V, n.~17); A (age allusion in \work{Ep.} 59 reported)}{Probable range; year unresolved}
\lifeevent{c.~353--354}{After father's death, mother brings Marcellina, Satyrus, and Ambrose to Rome; Pope Liberius veils Marcellina at St.~Peter's on Christmas Day}{A (\work{De Virginibus} III.1.1--4); K (\work{Prolegomena}~\S III)}{Event secure; exact year approximate}
\lifeevent{c.~365--370}{Advocate at the praetorian prefecture of Italy under Probus; appointed \latin{consularis} of Aemilia and Liguria at Milan}{A (\work{De Excessu} I.25, 49; \work{De Officiis} I.1.4); E (Paulinus via \work{Prolegomena}~\S V)}{Secure as career; exact years approximate}
\lifeevent{7~Dec.\ 374 (alt.\ 373)}{Death of Homoian bishop Auxentius; unbaptized governor acclaimed bishop by Milanese crowd, ratified by Valentinian~I, baptized Catholic, consecrated eight days later}{A (\work{Letter} XXI.7; \work{De Officiis} I.1.2--4); N (Basil, \work{Letter} CXCVII.1); M (\work{Martyrologium Romanum}, 7~Dec.)}{Event secure; 374 vs.\ 373 debated}
\lifeevent{374--375}{Sends clerics to St.~Basil in Cappadocia to recover the relics of exiled Catholic bishop St.~Dionysius of Milan}{N (St.~Basil, \work{Letter} CXCVII.2)}{Secure}
\lifeevent{c.~377}{\work{De Virginibus}, \work{De Viduis}; controversy with Milanese families over consecrated virginity}{A (\work{De Virginibus} I.11--12); K (\work{Prolegomena}~\S VI)}{Approximate}
\lifeevent{378}{First two books of \work{De Fide} for Emperor Gratian; Battle of Adrianople (9~Aug.); Ambrose melts sacred vessels to ransom captives}{A (\work{De Fide} I.prol.; \work{De Officiis} II.28.136--140)}{Secure}
\lifeevent{c.~378--379}{Death of brother Satyrus; two funeral orations \work{De Excessu Fratris Satyri}}{A (\work{De Excessu} I--II); M (\work{Martyrologium Romanum}, 17~Sept.); K (de Romestin: 17~Oct.\ 379)}{Event secure; exact date disputed}
\lifeevent{379--381}{Books~III--V of \work{De Fide}; \work{De Spiritu Sancto} (381); consecration of Anemius at Sirmium; Council of Aquileia condemns Palladius and Secundianus (Sept.\ 381)}{A (\work{De Fide}; \work{De Spiritu Sancto}); K (\work{Prolegomena}~\S III, \S VI)}{Secure}
\lifeevent{383--384}{Murder of Gratian (25~Aug.\ 383); first embassy to Maximus at Trier on behalf of Valentinian~II}{A (\work{Epistle} XX.23; \work{Letter} XXI.20)}{Secure}
\lifeevent{384}{Altar of Victory debate: Symmachus's \work{Relatio} III answered by \work{Epistles} XVII and XVIII; \work{De Paenitentia} against Novatians; Augustine arrives in Milan}{A (\work{Epistles} XVII--XVIII, LVII.2--3); N (Symmachus; Augustine, \work{Conf.} V.13.23)}{Secure}
\lifeevent{Holy Week 385}{First Basilica crisis: Justina demands Portian and New Basilicas; Ambrose refuses, rescues Arian priest Castulus; soldiers fraternize with Catholics}{A (\work{Epistle} XX)}{Secure}
\lifeevent{Jan.--June 386}{Homoian assembly law; \work{Letter} XXI refuses consistory trial against Auxentius/Mercurinus; \work{Sermon against Auxentius}; antiphonal hymns introduced; discovery of Sts.\ Gervasius and Protasius and healing of Severus (June)}{A (\work{Letter} XXI; \work{Sermon against Auxentius}; \work{Letter} XXII); N (Augustine, \work{Conf.} IX.7.15--16; \work{City of God} XXII.8)}{Secure}
\lifeevent{c.~386--387}{\work{Expositio Evangelii secundum Lucam}; \work{De Mysteriis}; Easter 387 baptism of Augustine, Alypius, and Adeodatus; second embassy to Maximus fails}{A (\work{Expositio Lucae}; \work{De Mysteriis}); N (Augustine, \work{Conf.} IX.5--6); K (\work{Prolegomena}~\S III)}{Secure; treatise dates approximate}
\lifeevent{Late 388}{Callinicum crisis: \work{Letters} XL and XLI; Ambrose halts Mass until Theodosius~I rescinds rebuilding order and inquiry}{A (\work{Letters} XL--XLI)}{Secure}
\lifeevent{c.~389--391}{\work{Hexameron} preached in Lent (c.~389); \work{De Officiis} (c.~391)}{A (\work{Hexameron}; \work{De Officiis}); K (\work{Prolegomena}~\S VI)}{Approximate}
\lifeevent{390}{Massacre at Thessalonica; confidential handwritten \work{Letter} LI excludes Theodosius~I from the Eucharist until public penance}{A (\work{Letter} LI); E (Paulinus; Theodoret V.18 at reported ceiling)}{Core secure; doorway tableau later}
\lifeevent{392--395}{Murder of Valentinian~II (392); Eugenius usurpation and Ambrose's withdrawal (\work{Letter} LVII, 393); Battle of the Frigidus and pleas for clemency (\work{Letters} LXI--LXII, 394); death of Theodosius~I (Jan.\ 395)}{A (\work{Letters} LVII, LXI--LXII); K (\work{Prolegomena}~\S III)}{Secure}
\lifeevent{396}{\work{Epistle} LXIII to the vacant see of Vercellae}{A (\work{Epistle} LXIII)}{Secure}
\lifeevent{4~Apr.\ 397}{Death on Holy Saturday after dictating Psalm~43; buried Easter Sunday, 5~April 397, beside Sts.\ Gervasius and Protasius in the Basilica Ambrosiana; succeeded by Simplicianus}{E (Paulinus via \work{Prolegomena}~\S III, \S V); N (Augustine, \work{Conf.} VIII.2.3 n.~608); M (\work{Martyrologium Romanum}, 4~Apr.)}{Secure}
\lifeevent{835 \& 1864}{Archbishop Angilbert~II places Ambrose and the two martyrs in a porphyry sarcophagus; 1864 excavations uncover sarcophagus and side vaults}{K/L (de Romestin \work{Prolegomena}~\S II, correcting its Cosmas/Damian slip)}{Reported only; no local archaeological report}
\lifeevent{1298}{Pope Boniface~VIII orders solemn feasts of the four original Latin Doctors of the Church}{M/K (reported reference consensus)}{Reported; decretal text not held locally}
\end{lifetimeline}

\subsection{Major cited writings of Ambrose}

\begin{worksinventory}
\workentry{\work{De Virginibus} \& \work{De Viduis}}{c.~377}{Addressed to Marcellina; defense of consecrated virginity and Christian widowhood}{Early episcopate}{I.2--3, 11--12; III.1.1--4}
\workentry{\work{De Excessu Fratris Satyri}}{c.~378--379}{Two funeral orations at Milan on the death of Satyrus and the Resurrection}{Early episcopate}{I.1--26, 31, 38, 43--49, 58--59; II.3}
\workentry{\work{De Fide} (Books~I--V)}{378--380}{Written at Emperor Gratian's request defending Nicene Christology}{Anti-Arian polemic}{I.prol.1--4; IV.10.124--125}
\workentry{\work{De Spiritu Sancto} (Books~I--III)}{381}{Written for Gratian proving the full divinity of the Holy Spirit}{Anti-Arian polemic}{I.prol. (also via Augustine)}
\workentry{\work{De Paenitentia} (Books~I--II)}{c.~384}{Against Novatianist rigorism; defense of priestly absolution and mercy}{Pastoral polemic}{I.1.1--4; I.2.5--9; I.3.10}
\workentry{\work{Epistles} XVII--XVIII}{384}{Replies to Symmachus's \work{Relatio} III on the Altar of Victory}{Imperial crisis}{XVII.10, 13; XVIII.7--8, 11--14, 16, 20--21}
\workentry{\work{Epistle} XX; \work{Letter} XXI; \work{Sermon against Auxentius}; \work{Letter} XXII}{385--386}{Basilica crisis with Justina and Valentinian~II; hymnody; discovery of Sts.\ Gervasius and Protasius}{Imperial crisis}{XX.2--8, 13, 20--28; XXI.2--20; \work{Sermon} 2, 18, 33--36; XXII.1--2, 12--22}
\workentry{\work{Expositio Evangelii secundum Lucam}}{c.~386--387}{Ten books of homilies on Luke's Gospel}{Biblical exegesis}{II.19, 26, 41--42 (CSEL \& Oxford \work{Catena})}
\workentry{\work{De Mysteriis} (and \work{De Sacramentis} IV)}{c.~387}{Easter-week mystagogical catechesis on Baptism, Confirmation, and Eucharist}{Sacramental catechesis}{\work{De Myst.} I--IX; \work{De Sacr.} IV.5.21--23, 6.27 (PL~16)}
\workentry{\work{Letters} XL--XLI}{388}{Callinicum synagogue dispute before Emperor Theodosius~I}{Imperial crisis}{XL.2--23; XLI.1, 26--28}
\workentry{\work{Hexameron} (Books~I--VI)}{c.~389}{Nine Lenten homilies on the six days of creation}{Biblical exegesis}{I.6.20; V.24.88--25.89; VI.2.7--8 (PL~14)}
\workentry{\work{Letter} LI}{390}{Private handwritten letter to Theodosius~I after the Thessalonica massacre}{Imperial crisis}{LI.2--15}
\workentry{\work{De Officiis} (Books~I--III)}{c.~391}{Christian recasting of Cicero on clerical virtues and charity}{Pastoral ethics}{I.1.1--4; II.28.136--141}
\workentry{\work{Letters} LVII, LXI--LXIII}{393--396}{Eugenius usurpation, clemency after Frigidus, and the vacant see of Vercellae}{Final years}{LVII.2--12; LXI.5, 7; LXII.3; LXIII.1--9, 68--70}
\workentry{Authenticated Hymns}{385--397}{Four externally verified iambic dimeter hymns (\work{Aeterne rerum}, \work{Deus creator}, \work{Iam surgit}, \work{Veni Redemptor})}{Liturgical poetry}{\work{Deus creator} in \work{Conf.} IX.12.32; \work{Hexameron} V.24.88}
\end{worksinventory}

\subsection{Received-tradition audit summary}

\begin{traditionaudit}
\traditionrow{Swarm of bees on infant's mouth}{Paulinus of Milan, \work{Vita} ch.~3 (c.~422; reported level)}{Absent from Ambrose, Augustine, and Basil; classical literary \latin{topos} (Plato, Pindar)}{Assigned to fifth-century hagiography and symbolic iconography}
\traditionrow{Child's cry ``Ambrose, bishop!'' and flight expedients}{Paulinus of Milan, \work{Vita} (reported via NPNF~II.10); 11th-c.\ \work{Manuale Ambrosianum}}{Ambrose (\work{Ep.} XXI.7; \work{De Off.} I.1.4) and Basil (\work{Ep.} CXCVII) confirm popular acclamation and reluctance without the child's shout}{Historical core in main narrative; Paulinus's dramatic particulars marked as Class~E}
\traditionrow{Alternating composition of \work{Te Deum} at Easter 387}{Early medieval legend (noted in NPNF~I.1 n.~756)}{Complete silence in Augustine's \work{Confessions} IX.6--7, 12; modern attribution to Nicetas of Remesiana (reported)}{Presented as liturgical legend; never used as biographical fact}
\traditionrow{Barring Theodosius at Milan cathedral doors (390)}{Theodoret, \work{Hist.\ Eccl.} V.18 (5th c., reported ceiling)}{Ambrose's contemporary handwritten \work{Letter} LI urges private withdrawal rather than staging a doorway clash}{Penance and exclusion affirmed from \work{Letter} LI; doorway scene flagged as 5th-c.\ dramatization}
\traditionrow{Sole authorship of the Ambrosian Rite and whole hymnal}{Popular medieval and modern attribution}{1640 \work{Missale Ambrosianum} preface credits cumulative Milanese pastors; only four core hymns carry 4th--5th-c.\ verification}{Ambrose's real liturgical core distinguished from centuries of Milanese evolution}
\traditionrow{Burial beside Gervasius \& Protasius; 835 sarcophagus; 1864 excavation}{Ambrose \work{Letter} XXII.13 (386); de Romestin \work{Prolegomena}~\S II (1896)}{De Romestin~\S II misnames the martyrs as Cosmas and Damian; no local archaeological excavation report}{397 burial beside Gervasius \& Protasius affirmed; 835/1864 chain kept at reported ceiling with slip corrected}
\end{traditionaudit}

\subsection{Rights and review status}

Original prose carries the CC~BY~4.0 license stated in the final colophon. Ancient patristic and liturgical texts are public domain; English translations follow the public-domain nineteenth-century NPNF printings (1886--1896) and the 1843 Oxford \work{Catena Aurea}, with direct translations from tracked public-domain Latin texts (Migne PL~14, CSEL~32/4, \work{Confessions} IX.12.32, \work{Manuale Ambrosianum}, \work{Missale Ambrosianum}, \work{Martyrologium Romanum}) expressly identified where given. Internal editorial checking, source-binding fingerprint verification, web conversion review, and print compilation checks have been completed; no external specialist, theological, canonical, archaeological, or ecclesiastical review is claimed.
